\section{Discussion}
\label{sec:discussion}

\subsection{Key Findings and Their Interpretation}

\textbf{Finding 1: Any tagging creates a substantial, robust discoverability advantage.} The 22.6\% adoption advantage (IRR=1.2263) is statistically overwhelming and survives the primary confound with only 12.2\% attenuation. The single most impactful metadata action for a dataset creator is simply adding at least one tag --- the choice of tags matters less than the act of tagging itself, because the primary mechanism is binary search index membership.

\textbf{Finding 2: Tag count magnitude is a genuine dose-response mechanism.} Within the tagged subset, IRR=1.5332 per log-unit with attenuation ratio 1.0007 establishes that additional tags expand search pathways independently of platform era. Each additional tag is an additional keyword query pathway, reaching different subsets of researchers. This is the amplification arm of the mechanism.

\textbf{Finding 3: Comprehensive tagging (6+) is the dominant amplification tier.} The 6+ tier produces a 28.6\% adoption advantage, meaningfully above the 12.7--12.8\% for 1--5 tags. The bimodal tagging distribution (all-or-nothing) means that minimal tagging (1--2 tags) produces effects statistically indistinguishable from moderate tagging (3--5 tags). Repository interfaces that encourage minimal tag entry may not achieve the full discoverability benefit; a threshold near 6 tags appears to be the meaningful target.

\subsection{Connection to Literature}

Our binary threshold finding provides the first NB-2 IRR for keyword tagging on an ML repository, filling the quantification gap in the \citet{Wilkinson2016FAIR} FAIR framework. The direction is consistent with \citet{Yang2024Navigating} and \citet{Lachmuth2025BonaRes}, though platform differences preclude direct comparison. Our mechanism verification via decade fixed effects (H-M1) provides a methodological template for observational platform studies: testing extreme temporal collinearity without abandoning FE controls demonstrates that within-cohort variation can identify structural mechanisms even under severe between-cohort confounding.

\subsection{Limitations}

\textbf{L1: Cross-sectional --- predictive, not causal without temporal data.} Tags and task counts are observed simultaneously. The causal interpretation rests on OpenML's creator-only tagging architecture \citep{Vanschoren2014OpenML} as a structural temporal ordering defense, unverified by timestamps in the available corpus.

\textbf{L2: \texttt{N\_tasks} proxy measures registration breadth, not execution depth.} Task registration is a valid, deliberate adoption signal --- but may diverge from execution-frequency adoption for some datasets.

\textbf{L3: Conditional adoption --- hurdle component deferred.} We study $N\_tasks \geq 1$ datasets (adoption intensity); the probability of any adoption at all is future work.

\textbf{L4: 12.2\% decade attenuation partially uncontrolled.} The residual after FE may include both structural FAIR F1 mechanism and uncontrolled era confounding. The gate passed with large margin despite this.

\textbf{L5: H-M3 bin 1--2 sparsity (N=73) limits uniform categorical dose-response claim.} The low-count range is undetermined; the 6+ tier characterization is robust.

\textbf{L6: For the continuous measure (H-M2, H-M3), reverse causality remains possible.} Popular datasets attracting additional tags retroactively --- requires tag assignment timestamps unavailable in the current corpus. OpenML's creator-only tagging architecture partially mitigates this, but if community tag additions exist, they could introduce upward bias in \texttt{log\_tag\_count\_p1} estimates.

\subsection{Broader Impact}

\textbf{Positive impacts:} IRR estimates translate FAIR F1 compliance from aspirational to actionable. Repository designers can justify tag-requirement policies; dataset creators gain evidence-based guidance (tag, and target 6+). The NB-2 methodology with decade FE provides a replication template for other repositories.

\textbf{Potential concerns:} Strategic tag farming analogous to SEO manipulation could emerge if keyword tagging becomes widely known to boost adoption. Repository governance should monitor for tag quality degradation alongside quantity. Additionally, if discovery concentrates among heavily-tagged datasets, a Matthew effect could emerge --- popular datasets becoming more discoverable at the expense of quality untagged datasets worth monitoring.
